% Generated draft: requires mathematical and source review.
\documentclass[11pt]{article}
\usepackage[margin=25mm]{geometry}
\usepackage{fontspec}
\setmainfont{Latin Modern Roman}
\usepackage{amsmath,amssymb,mathtools,unicode-math}
\setmathfont{Latin Modern Math}
\usepackage{enumitem,needspace,xurl,hyperref}
\setlength{\parindent}{0pt}
\setlength{\parskip}{5pt}
\setlength{\emergencystretch}{3em}
\newcommand{\N}{\mathbb{N}}
\newcommand{\Z}{\mathbb{Z}}
\newcommand{\Q}{\mathbb{Q}}
\newcommand{\R}{\mathbb{R}}
\newcommand{\C}{\mathbb{C}}
\newcommand{\F}{\mathbb{F}}
\newcommand{\A}{\mathbb{A}}
\newcommand{\PP}{\mathbb{P}}
\newcommand{\E}{\mathbb{E}}
\newcommand{\eps}{\varepsilon}
\newcommand{\dd}{\,\mathrm{d}}
\begin{document}
\section*{UnsolvedMath ID: 20001936}
\textbf{Source problem code:} catalogue entry
\subsection{Definitions and mathematical statement}
Question \(12.0.3\) asks whether a closed pseudo-Riemannian manifold \((M,g)\) with a nonzero parallel self-adjoint nilpotent endomorphism \(A:TM\to TM\) must admit a nonzero parallel lightlike vector field, possibly after a double cover. Explicitly,
\[
 \nabla A=0,\qquad g(Au,v)=g(u,Av),\qquad A^j=0\text{ for some }j\geq2.
\]
A requested vector field \(V\) satisfies \(\nabla V=0\), \(V\not\equiv0\), and \(g(V,V)=0\). Positive-definite metrics admit no nonzero self-adjoint nilpotent, so the meaningful setting is indefinite; the source gives no signature restriction.

The aid reports a general compact flat counterexample mechanism. Let \(E/\Gamma\) be a compact flat Riemannian manifold, with translation lattice \(\Lambda\subset E\), finite orthogonal linear holonomy \(H\), odd \(|H|\), and \(E^H=0\). If \(\gamma x=A_\gamma x+b_\gamma\), act on \(E\oplus E\) by
\[
 (x,y)\longmapsto(A_\gamma x+b_\gamma,\ A_\gamma y+\lambda),
 \qquad \gamma\in\Gamma,\ \lambda\in\Lambda.
\]
The resulting quotient is closed and complete flat. For \(m=\dim E\), it carries the invariant neutral metric and endomorphism
\[
 g((u,v),(u',v'))=\langle u,v'\rangle+\langle v,u'\rangle,
 \qquad A(u,v)=(v,0).
\]
They have signature \((m,m)\), \(A\ne0\), \(A^2=0\), and \(A\) is self-adjoint and parallel. The linear holonomy is diagonal \(H\) on \(E\oplus E\), with no fixed vector; hence there is no nonzero parallel vector field at all. For any connected double cover the holonomy image has index dividing two in the odd-order group \(H\), so remains all of \(H\); the cover also has no such field.

The aid instantiates this construction using Bauer--Gleissner's reported compact flat real eight-manifold with holonomy \((\Z/3\Z)^2\) and no invariant vector. It therefore gives a closed complete flat signature-\((8,8)\) counterexample, including every connected double cover. A full reference and verification of the required eight-dimensional holonomy representation are not supplied.

The aid also gives a positive rank-one case. If \(A^{j-1}\ne0\) is its highest nonzero power and \(\operatorname{rank}\operatorname{im}A^{j-1}=1\), that image is a parallel null line whose holonomy acts only by signs. Its line-orientation double cover, possibly trivial, admits a nonzero parallel null vector. In Lorentzian signature the top image is totally isotropic and hence has rank one, so the source question is affirmative in that restricted signature. This does not contradict the neutral higher-index example.

The appended Finsler paragraph defines affine equivalence by equality of affinely parametrized, constant-speed geodesics and comments that parametrized geodesics determine the relevant connection. In the general Finsler setting this is naturally understood at the geodesic-spray level, not as an assertion that every Finsler metric has a Riemannian Levi--Civita connection. It is neighboring context rather than a further hypothesis of the nilpotent question.
\subsection{Short English statement}
The aid gives a neutral-signature flat counterexample with a parallel self-adjoint square-zero endomorphism but no parallel vectors, even on a connected double cover. Odd holonomy ensures the obstruction survives that cover. A rank-one top image instead yields a null vector on its line-orientation cover, including the Lorentzian case.
\subsection{Sources}
\begin{itemize}
\item Local UnsolvedMath registry, permanent ID \(20001936\), source code \texttt{catalogue entry}; canonical collection \texttt{aim-geometry-notes.json}, notes[273]. Source item \(12.0\), actual question \(12.0.3\).
\item Keith Burns and Vladimir Matveev, ``Open problems and questions about geodesics,'' AIM workshop ``Geodesics'': \url{https://aimath.org/pastworkshops/geodesicsproblems.pdf}. The supplied verification reports PDF inspection; no PDF was opened in this drafting session. The following affine-Finsler paragraph is retained as context.
\item Bauer--Gleissner's compact flat real eight-manifold with holonomy \((\Z/3\Z)^2\) is named by the aid without a full title, identifier or publication citation. Original remarks and literature fields are empty; no proposer or year is supplied.
\item Supplied research attempt \(1\), result ``counterexample,'' run status ``refuted,'' rigor label ``proved,'' novelty confidence ``medium.'' Registry status ``solved,'' classification \texttt{SOLVED-IN-LITERATURE}, updated 14 August 2026. These labels are supplied assessments, not independent verification of the proposed construction.
\end{itemize}
\end{document}
